\section{Open questions}

The following five questions are genuinely open after this replication.
Each has an experimental protocol that fits within the standing
free-endpoints-only constraint (Qiskit Aer, Mitiq, Cirq, and IBM Open
Plan free hardware only).

\begin{enumerate}
  \item \textbf{Noise-model robustness.}  Does the paper's
    method-ranking-by-shot-budget claim (VD/ZNE win at
    $N_{\rm tot}=10^5\text{--}10^7$; CDR/vnCDR at $10^7\text{--}10^9$;
    UNITED at $10^{10}$) survive when the underlying noise model is
    swapped from IonQ trapped-ion to a superconducting-transmon noise
    model (IBM Falcon/Eagle-style with correlated $ZZ$ crosstalk,
    $T_1/T_2$ asymmetry, and measurement error)?
    \emph{Basis.} Paper's benchmark is single-noise-model; QEM
    method rankings are famously noise-structure-sensitive
    (PEC's failure in our stack is an example).
    \emph{Next step.} Reproduce the ZNE/CDR/vnCDR/UNITED pipeline on
    Aer \texttt{FakeSherbrooke}/\texttt{FakeCairoV2} noise (free,
    calibrated to real IBM hardware) and the paper's trapped-ion
    model as anchor; diff the winning-method table.

  \item \textbf{Extension to newer QEM techniques.}  Does modern QEM
    beyond the paper's set --- Subspace Expansion, Symmetry
    Verification, learning-based PEC with GST-learned representations
    --- change the accuracy-vs-cost Pareto frontier the paper reports?
    \emph{Basis.} Paper predates or omits several QEM techniques that
    have matured since (Mitiq's learning-based PEC and
    subspace-expansion modules).
    \emph{Next step.} Extend the paper's harness with these methods on
    the same random-quantum-circuit and QAOA-Max-Cut suites at
    $Q\in\{4,6,8\}$; regenerate the paper's shot-budget curves with
    the enlarged method set.

  \item \textbf{Problem-adaptive selection.}  Is there a
    problem-adaptive rule --- a classical decision tree taking
    $(Q, {\rm depth}, {\rm circuit\text{-}family\ fingerprint},
    N_{\rm tot})$ as input and outputting the predicted best-of-set
    QEM method --- that systematically beats ``always pick UNITED''?
    \emph{Basis.} Paper shows the winner depends on both shot budget
    and problem class but does not propose an adaptive selector.
    \emph{Next step.} Compute the empirical winner per tuple on the
    paper's benchmark, train a small classifier on cheap circuit
    features (2-design distance, Clifford fraction, ideal-state
    entanglement entropy, 2-qubit gate density), and measure
    hold-out $|{\rm err}|$ gain of adaptive selection vs.\
    always-UNITED.

  \item \textbf{Integration with algorithmic EM.}  How does UNITED
    interact with algorithmic error mitigation (echo verification,
    mid-circuit-measurement post-selection, code-based error
    detection) --- is there a stacking rule where
    ``algorithmic-inside, UNITED-outside'' gives an additive-in-dB
    error reduction, or does it saturate?
    \emph{Basis.} The next 2 years of QEM literature has emphasized
    stacked mitigations for early-fault-tolerant quantum computing;
    the paper does not address composability.
    \emph{Next step.} Stacked benchmark: (i) algorithmic EM alone,
    (ii) UNITED alone, (iii) algo+UNITED stacked; check whether
    stacked $\geq \max({\rm algo},{\rm UNITED})$ and whether the
    error-suppression factor multiplies.

  \item \textbf{Hardware runtime on modern NISQ.}  Does the paper's
    method-ranking hold on 2026 NISQ hardware (IBM Heron~r2/Fez), where
    real device drift, non-Markovian noise, and calibration-cycle
    timing dominate?  Specifically, is CDR still viable when the
    near-Clifford classical training oracle has to model calibration
    drift within a single execution window?
    \emph{Basis.} Paper is a simulator-based benchmark; it silently
    assumes stationary Markovian noise per shot, which is optimistic
    on real hardware.
    \emph{Next step.} Use free IBM Open Plan minutes to run a
    scaled-down $Q=4$, depth $\leq 16$ subset with ZNE and CDR under
    Mitiq; log \texttt{ibm\_calibrations} JSON at run-start and
    run-end; compare hardware improvement-factor vs.\ simulator
    prediction; report whether the multiplier degrades ($\sim 2\times$)
    and whether CDR's classical-oracle assumption survives observed
    drift.
\end{enumerate}
